\chapter{Wyrazenia i operatory}

\section{Czym jest wyrazenie}
Wyrazenie to fragment kodu, ktory po obliczeniu daje jakas wartosc. Na przyklad:

\begin{itemize}
  \item \mc{2 + 3} daje liczbe,
  \item \mc{title + "!"} daje tekst,
  \item \mc{views > 100} daje wartosc logiczna,
  \item \mc{isPublished and isFeatured} daje wartosc logiczna,
  \item \mc{condition ? a : b} wybiera jedna z dwoch wartosci.
\end{itemize}

W MoonChunk wyrazenia pojawiaja sie praktycznie wszedzie: po prawej stronie przypisania,
w warunkach, w petlach, w argumentach funkcji, w atrybutach HTML i we wstawieniach do
\mc{content}.

\section{Operatory arytmetyczne}
MoonChunk obsluguje podstawowe operacje liczbowe:

\begin{itemize}
  \item \mc{+} dodawanie,
  \item \mc{-} odejmowanie,
  \item \mc{*} mnozenie,
  \item \mc{/} dzielenie,
  \item \mc{\%} reszta z dzielenia.
\end{itemize}

\begin{lstlisting}[style=moonchunk,caption={Podstawowe obliczenia}]
let sum: int = 2 + 3;
let diff: int = 10 - 4;
let mul: int = 6 * 7;
let div: int = 8 / 2;
let mod: int = 7 % 2;
\end{lstlisting}

\section{Kolejnosc wykonywania dzialan}
Podobnie jak w matematyce, nie wszystkie operacje sa wykonywane od lewej do prawej w
tej samej kolejnosci. Najpierw interpreter bierze pod uwage priorytety operatorow.
Najprostsza intuicja jest taka:

\begin{itemize}
  \item nawiasy maja najwyzszy priorytet,
  \item potem mnozenie, dzielenie i modulo,
  \item potem dodawanie i odejmowanie,
  \item dalej porownania,
  \item a na koncu bardziej zlozone konstrukcje logiczne i warunkowe.
\end{itemize}

Jesli chcesz uniknac watpliwosci, uzywaj nawiasow.

\begin{lstlisting}[style=moonchunk,caption={Nawiasy zwiekszaja czytelnosc}]
let value: int = (2 + 3) * 4;
\end{lstlisting}

\section{Dzielenie przez zero i modulo przez zero}
MoonChunk traktuje te sytuacje jako blad. To dobre, bo wynik typu \emph{nieskonczonosc}
bylby w praktyce niebezpieczny i moglby ukryc pomylke w logice programu.

\begin{lstlisting}[style=moonchunk,caption={Przyklady niedozwolonych operacji}]
let a: int = 10 / 0;
let b: int = 10 % 0;
\end{lstlisting}

Takie zapisy powinny zakonczyc sie komunikatem diagnostycznym.

\section{Operatory porownania}
Porownania sa potrzebne wtedy, gdy chcemy sprawdzic relacje miedzy wartosciami.
MoonChunk obsluguje:

\begin{itemize}
  \item \mc{==} rownosc,
  \item \mc{!=} nierownosc,
  \item \mc{<} mniejsze niz,
  \item \mc{>} wieksze niz,
  \item \mc{<=} mniejsze lub rowne,
  \item \mc{>=} wieksze lub rowne.
\end{itemize}

\begin{lstlisting}[style=moonchunk,caption={Porownania w praktyce}]
let isZero: bool = count == 0;
let isBig: bool = views > 100;
let isReady: bool = step >= 3;
\end{lstlisting}

Kazde takie porownanie daje wynik typu \mc{bool}.

\section{Operatory logiczne}
Do skladania warunkow MoonChunk uzywa operatorow tekstowych:

\begin{itemize}
  \item \mc{and},
  \item \mc{or},
  \item \mc{not}.
\end{itemize}

\begin{lstlisting}[style=moonchunk,caption={Logika w MoonChunk}]
let visible: bool = isPublished and not isArchived;
let premium: bool = hasSubscription or hasInvitation;
\end{lstlisting}

To podejscie jest czytelne i dobrze pasuje do skladni zorientowanej na czlowieka.

\section{Operator warunkowy}
MoonChunk wspiera operator trójargumentowy, znany z wielu innych jezykow.
Pozwala on wybrac jedna z dwoch wartosci na podstawie warunku.

\begin{lstlisting}[style=moonchunk,caption={Operator warunkowy}]
let label: string = isPublished ? "Opublikowany" : "Szkic";
\end{lstlisting}

To zwięzly zapis przydatny wtedy, gdy chcesz szybko wybrac jedna z dwoch prostych wartosci.
Jesli logika robi sie bardziej rozbudowana, zwykle lepiej przejsc do pelnego \mc{if}.

\section{Laczenie tekstu}
Operator \mc{+} sluzy nie tylko do dodawania liczb. W praktyce jest tez wygodnym sposobem
na skladanie tekstow.

\begin{lstlisting}[style=moonchunk,caption={Laczenie tekstu}]
let fullTitle: string = "MoonChunk" + " " + "Docs";
\end{lstlisting}

Z takiego mechanizmu czesto korzysta sie przy budowaniu linkow, naglowkow, klas CSS
czy fragmentow HTML skladanych do jednej zmiennej tekstowej.

\section{Przypisanie jako wyrazenie i zwykle przypisanie}
MoonChunk pozwala przypisywac nowa wartosc do juz istniejacej zmiennej:

\begin{lstlisting}[style=moonchunk,caption={Zwykle przypisanie}]
let counter: int = 0;
counter = counter + 1;
\end{lstlisting}

To nie jest redeklaracja, tylko zmiana wartosci istniejacej zmiennej typu \mc{let}.
Roznica jest fundamentalna:

\begin{itemize}
  \item \mc{let counter = 0;} tworzy nowa nazwe,
  \item \mc{counter = counter + 1;} zmienia wartosc tej nazwy.
\end{itemize}

\section{Operator inkrementacji}
MoonChunk wspiera dwa znane zapisy zwiekszania wartosci o jeden:

\begin{itemize}
  \item \mc{i++} -- postinkrementacja,
  \item \mc{++i} -- preinkrementacja.
\end{itemize}

\begin{lstlisting}[style=moonchunk,caption={Preinkrementacja i postinkrementacja}]
let i: int = 1;
print(i++);
print(++i);
\end{lstlisting}

W praktyce:

\begin{itemize}
  \item \mc{i++} zwraca stara wartosc, a dopiero potem zwieksza zmienna,
  \item \mc{++i} najpierw zwieksza zmienna, a potem zwraca nowa wartosc.
\end{itemize}

To drobna roznica, ale przy bardziej zaawansowanej logice bywa istotna.

\section{Wyrazenia w HTML}
W MoonChunk wyrazenia moga trafiać bezposrednio do \mc{content}. Przykład:

\begin{lstlisting}[style=moonchunk,caption={Wyrazenie wewnatrz content}]
content {
  <p>{views > 100 ? "Popularny wpis" : "Nowy wpis"}</p>
};
\end{lstlisting}

To bardzo wygodne, ale warto zachowac umiar. Jesli wyrazenie staje sie dlugie,
czesto lepiej obliczyc je najpierw w zmiennej.

\section{Wyrazenia w indeksowaniu i dostepie do pol}
Wyrazenia moga pojawiac sie rowniez przy dostepie do tablic i obiektow.

\begin{lstlisting}[style=moonchunk,caption={Dostep dynamiczny do danych}]
const cards = [
  { title: "Parser" },
  { title: "Runtime" }
];

let index: int = 1;
let key: string = "title";
print(cards[index][key]);
\end{lstlisting}

To pokazuje, ze MoonChunk nie ogranicza sie do prostych stalych sciezek typu
\mc{obj.title}. Potrafi tez pracowac z kluczami i indeksami obliczanymi dynamicznie.

\section{Kiedy wyrazenie jest zbyt zlozone}
Jesli podczas czytania wlasnego kodu sam musisz zatrzymac sie na kilka sekund, aby
zrozumiec jedno wyrazenie, to sygnal, ze warto je uproscic. Dobra praktyka wyglada tak:

\begin{lstlisting}[style=moonchunk,caption={Uproszczenie przez zmienne pomocnicze}]
let canShowBanner: bool = isPublished and not isArchived;
let bannerText: string = canShowBanner ? "Gotowe" : "Ukryte";
\end{lstlisting}

Taki zapis jest zwykle czytelniejszy niz jeden bardzo dlugi warunek wpisany wprost do
HTML-a albo do warunku petli.
